\documentclass[11pt,a4paper]{article}
\usepackage{amsmath,amssymb,amsthm}
\usepackage{geometry}
\geometry{margin=1in}
\usepackage{hyperref}
\usepackage{enumitem}

\newtheorem{theorem}{Theorem}[section]
\newtheorem{lemma}[theorem]{Lemma}
\newtheorem{conjecture}[theorem]{Conjecture}
\newtheorem{problem}[theorem]{Problem}
\theoremstyle{remark}
\newtheorem{remark}[theorem]{Remark}

\title{\bfseries A Greedy Dirichlet Sub‑Sum Transform Approach\\
to the Riemann Hypothesis}
\author{}
\date{}

\begin{document}
\maketitle

\begin{abstract}
The harmonic–categorical framework, through its Harmonic Sine Transform (HST), provides a transfer operator \(\mathcal{L}_s\) whose Fredholm determinant equals \(\zeta(s)/\zeta(2s)\) up to an entire zero‑free factor.  The failure of the finite‑dimensional spectral correspondence (Programme~VII) blocked the original route to a proof of the Riemann Hypothesis.  Here we introduce a new perspective: the \emph{Greedy Dirichlet Sub‑Sum Transform (GDST)}, which recasts the HST as a family of Dirichlet series \(E(\theta,s)\) built from the greedy digits \(\delta_n(\theta)\).  This formulation links the geometric expansion of angles directly to the Dirichlet eta function and opens a fresh pathway toward a self‑adjoint spectral realisation.  We outline a research programme whose aim is to prove that the zeros of the Riemann zeta function coincide with the eigenvalues of a canonically constructed operator derived from the GDST.
\end{abstract}

\section{Introduction}
The original ten‑step programme (Programmes~I–X) collapsed at Step 7, where the eigenvalues of the truncated Hermitian matrices \(H_k\) could not be identified with the zeros of the truncated determinants.  All attempts to mend that gap—either by approximating the full unitary group or by importing external Hamiltonians—failed.  Consequently, the Riemann Hypothesis remains open.

The present document re‑examines the foundations and discovers a more direct analytic object: the \emph{Greedy Dirichlet Sub‑Sum Transform} \(E(\theta,s)\).  For each angle \(\theta\in[0,\pi]\) the greedy digits select a sub‑sum of the shifted zeta (or eta) series, giving a one‑parameter family of Dirichlet series whose meromorphic continuations have poles precisely at the non‑trivial zeros of \(\zeta(s)\).  This object lives entirely in the space of Dirichlet series and eliminates the need for the problematic finite‑dimensional truncations.  We propose a new programme to extract a self‑adjoint operator from the GDST and to prove that its spectrum equals the set of ordinates \(\gamma\) with \(\zeta(\frac12+i\gamma)=0\).

\section{The Greedy Dirichlet Sub‑Sum Transform (GDST)}
Let \(\delta_n(\theta)\in\{0,1\}\) be the greedy digits defined by the harmonic fractions \(\pi/(n+2)\) as in \cite{Geere2026a}.  For \(\operatorname{Re}s>1\) define
\begin{equation}\label{E}
E(\theta,s) \;=\; \sum_{n=0}^{\infty} \frac{\delta_n(\theta)}{(n+2)^s}
        \;=\; \sum_{\{n:\,\delta_n(\theta)=1\}} \frac{1}{(n+2)^s}.
\end{equation}
The signed version, convergent for \(\operatorname{Re}s>0\), is
\begin{equation}\label{Epm}
E_{\pm}(\theta,s) \;=\; \sum_{n=0}^{\infty} (-1)^n \frac{\delta_n(\theta)}{(n+2)^s}.
\end{equation}
The family \(\{E(\theta,s)\}_{\theta\in[0,\pi]}\) (or its signed counterpart) constitutes the \textbf{Greedy Dirichlet Sub‑Sum Transform (GDST)} of the angle \(\theta\).

\begin{remark}
The map \(\theta\mapsto E(\theta,s)\) is a step function increasing at the threshold angles \(\theta_n^*\) with jumps \((n+2)^{-s}\).  Its total variation up to \(\theta=\pi\) equals \(\sum_{n=0}^{\infty} (n+2)^{-s} = \zeta(s)-1\).  Thus the GDST interpolates between the empty sum (at \(\theta=0\)) and the full shifted zeta series (at \(\theta=\pi\)).
\end{remark}

\section{Connection to the Riemann Zeta Function}
The transfer operator \(\mathcal{L}_s\) and the GDST are linked by the resolvent identity
\[
E(\theta,s) \;=\; \frac{\theta/\pi}{s-1} \;+\; \bigl\langle (I-\mathcal{L}_s)^{-1}\mathcal{L}_s\,\mathbf{1}\bigr\rangle (y_0),
\]
where \(y_0=2\theta/\pi\) \cite{HSTpaper}.  Hence the meromorphic continuation of \(E(\theta,s)\) mirrors that of \((I-\mathcal{L}_s)^{-1}\): its poles occur exactly at the non‑trivial zeros of \(\zeta(s)\) (and possibly at the trivial zeros, which are cancelled by the denominator \(\zeta(2s)\) in the determinant).  In particular, for each non‑trivial zero \(\rho\) of \(\zeta(s)\), the function \(\theta\mapsto \operatorname{Res}_{s=\rho}E(\theta,s)\) is non‑trivial and encodes information about the zero.

The \textbf{central observation} is that the GDST provides a \emph{scalar} family of Dirichlet series that already contains the full zero‑set of \(\zeta(s)\).  One may therefore attempt to construct a spectral problem directly on the angle \(\theta\), whose eigenfunctions are the indicator functions of the cylinder intervals, rather than on an auxiliary operator space.

\section{Research Programme: From GDST to a Self‑Adjoint Operator}
The programme consists of six problems, each building on the previous ones.

\subsection{Problem G1: Analytic structure of the GDST}
Give a complete, rigorous proof that for every \(\theta\in[0,\pi]\) the function \(s\mapsto E(\theta,s)\) extends meromorphically to the entire complex plane, with a simple pole at \(s=1\) of residue \(\theta/\pi\), and all other poles located exactly at the non‑trivial zeros of \(\zeta(s)\).  Determine the exact poles and residues of the signed version \(E_\pm(\theta,s)\) in the critical strip \(0<\operatorname{Re}s<1\).

\subsection{Problem G2: The balance angle mapping}
For a non‑trivial zero \(\rho\) of \(\zeta(s)\), define the \emph{balance angle} \(\theta_{\mathrm{bal}}(\rho)\in[0,\pi]\) as the unique solution of
\[
\operatorname{Re} E_\pm(\theta,\rho) \;=\; \frac12 .
\]
(Existence and uniqueness are suggested by numerical experiments \cite{Geere2026GHeta}.)  Study the mapping \(\rho\mapsto\theta_{\mathrm{bal}}(\rho)\).  Is it injective?  Does it order the zeros?  Prove that the set of balance angles is dense in \([0,\pi]\) if and only if RH holds.

\subsection{Problem G3: The GDST distribution as a spectral shift function}
Consider the increasing step function \(N(t;\theta)=\#\{n\ge 0:\theta_n^*\le\theta,\; n+2\le e^{t}\}\).  Its asymptotic behaviour is
\[
N(t;\theta) \sim \frac{\theta}{\pi}\,t \quad (t\to\infty).
\]
Define the \emph{GDST phase function}
\[
\phi(t;\theta) = \arg \bigl( E(\theta,\tfrac12+it) \bigr),\qquad t\in\mathbb{R},
\]
with a suitable branch continuous away from the zeros.  Show that the distributional derivative \(\frac{d}{dt}\phi(t;\pi)\) is a sum of delta functions supported at the ordinates \(\gamma\) of the non‑trivial zeros, weighted by their multiplicities.  This would directly identify the zero‑counting measure with the jump spectrum of the phase of the GDST at \(\theta=\pi\).

\subsection{Problem G4: A semi‑infinite Jacobi matrix realisation}
The greedy digits \(\delta_n(\theta)\) define a binary tree.  The associated Haar wavelets provide an orthonormal basis of \(L^2[0,\pi]\).  The multiplicative operator by the independent variable \(\theta\) in this basis yields a Jacobi (tridiagonal) matrix.  Prove that the spectrum of this Jacobi matrix is purely discrete and consists exactly of the balance angles \(\theta_{\mathrm{bal}}(\rho)\).  Then the Riemann zeros would be encoded in the spectral measure of a simple Jacobi operator.

\subsection{Problem G5: A Dirichlet‑to‑Neumann map for the greedy tree}
View the greedy tree as a metric graph with edge lengths given by the harmonic fractions.  On the boundary \([0,\pi]\) define a Dirichlet‑to‑Neumann map associated with a discretised Laplacian on the tree.  Show that the eigenvalues of this map are the ordinates of the Riemann zeros.  This would realise the Connes–Consani programme within the harmonic‑categorical framework.

\subsection{Problem G6: Proof of the Riemann Hypothesis}
Using the spectral realisation obtained in Problem G4 or G5, prove that the operator in question is self‑adjoint, hence its spectrum is real.  Consequently, all non‑trivial zeros of \(\zeta(s)\) lie on the critical line \(\operatorname{Re}s=\frac12\).

\section{Discussion and Outlook}
The GDST programme differs fundamentally from the earlier finite‑dimensional truncation approach.  Instead of approximating an infinite‑dimensional transfer operator by matrices, we work directly with the angle‑parametrised Dirichlet series \(E(\theta,s)\).  The key idea is that the parameter \(\theta\) itself becomes the spectral variable of a self‑adjoint operator—either a Jacobi matrix or a Dirichlet‑to‑Neumann map—whose eigenvalues are the balance angles.  The Riemann zeros then appear as the solutions of \(\operatorname{Re}E_\pm(\theta,\frac12+it)=\frac12\), a condition that reduces the two‑parameter family to a one‑dimensional set.

At present, all steps beyond G1 are conjectural.  However, the GDST formulation has the advantage of being both elementary and deeply connected to the arithmetic of the zeta function.  Numerical exploration of the balance angle mapping and the Jacobi matrix spectrum will provide crucial guidance.

\section*{Acknowledgements}
The author is indebted to Victor Geere for sharing the original notes on the greedy harmonic decomposition and for the insights into the Dirichlet eta function connection.

\begin{thebibliography}{9}
\bibitem{Geere2026a}
V.~Geere, \emph{On the Correlation Structure of a Greedy Harmonic Decomposition of the Sine Function}, March 2026.
\bibitem{Geere2026b}
V.~Geere, \emph{A Harmonic Reconstruction of the Sine Function and Its Relation to the Riemann Zeta Function}, 2026.
\bibitem{Geere2026GHeta}
V.~Geere, \emph{The Greedy Harmonic Decomposition and the Dirichlet Eta Function}, March 2026.
\bibitem{HSTpaper}
V.~Geere et al., \emph{A Harmonic Sine Transform and the Riemann Zeta Function}, 2026.
\end{thebibliography}

\end{document}